% TOP_PROBLEM: {"problemId": "problem.top500-omission-w045049-dickson-linear-prime-values", "canonicalTitle": "Dickson's Conjecture on Simultaneous Prime Values of Linear Forms", "releaseRank": 62, "primaryDomain": "number_theory_arithmetic_geometry", "primaryDomainLabel": "Number theory & arithmetic geometry", "displayStatus": "open_with_solved_subcases", "releaseStatus": "open_with_solved_subcases"}
% !TeX program = lualatex
% Definition and sources only; this file is not an LLM solution attempt.
\documentclass[11pt]{article}
\usepackage[margin=25mm]{geometry}
\usepackage{fontspec}
\setmainfont{Latin Modern Roman}
\usepackage{amsmath,amssymb,mathtools}
\usepackage{unicode-math}
\setmathfont{Latin Modern Math}
\usepackage{xurl,hyperref}
\hypersetup{colorlinks=true,urlcolor=blue}
\setcounter{secnumdepth}{0}
\setlength{\parindent}{0pt}
\setlength{\parskip}{5pt}
\setlength{\emergencystretch}{3em}
\begin{document}
\section*{\texorpdfstring{Dickson's Conjecture on Simultaneous Prime Values of Linear Forms}{Dickson's Conjecture on Simultaneous Prime Values of Linear Forms}}\label{62}
\subsection{Definitions and mathematical statement}
For $f_i(n)=a_in+b_i$ with integers $a_i\ge1$, define admissibility by the absence of a fixed prime divisor of $\prod_i f_i(n)$. Explicitly,
\[
 \forall p\in\mathcal P\ \exists n\in{\mathbb{Z}}:\quad p\nmid\prod_i(a_in+b_i).
\]
Dickson's assertion is $\forall X\ \exists n>X$ with all $a_in+b_i$ prime. The number of forms is finite and fixed for each assertion. Positivity for large $n$ follows from the positive slopes. The record does not require distinct forms or an asymptotic counting formula.
\par\smallskip\textit{Formulation and concept references:} \href{https://arxiv.org/html/0906.3850}{[S1]}.

\subsection{Short English statement}
Must every locally admissible finite family of positive-slope linear forms produce simultaneous primes arbitrarily often?
\subsection{Sources}
\begin{itemize}
\item[S1]Shaohua Zhang, Notes on Dickson's Conjecture, arXiv:0906.3850v3 (2009).
\url{https://arxiv.org/html/0906.3850}.
\textit{Formulation source.}
\item[S2]An exposition on the supersimplicity of certain expansions of the additive group of the integers.
\url{https://arxiv.org/html/2502.03915}.
\textit{Further reference; not a proof endorsement.}
\item[S3]From Golomb to Schinzel.
\url{https://www.preprints.org/manuscript/202505.0651}.
\textit{Further reference; not a proof endorsement.}
\end{itemize}
\end{document}
